\documentclass[border=10pt]{standalone}
\usepackage{tikz}
\usetikzlibrary{shapes, arrows, positioning, fit, calc, shadows}

\tikzset{
    chapter/.style={
        rectangle,
        draw=blue!60,
        fill=blue!10,
        thick,
        minimum width=3cm,
        minimum height=1.2cm,
        rounded corners,
        align=center,
        font=\small\bfseries,
        drop shadow
    },
    section/.style={
        rectangle,
        draw=green!60,
        fill=green!5,
        thick,
        minimum width=2.5cm,
        minimum height=0.8cm,
        rounded corners,
        align=center,
        font=\footnotesize,
        drop shadow
    },
    arrow/.style={
        ->,
        >=stealth,
        thick,
        color=blue!70
    }
}

\begin{document}
\begin{tikzpicture}[node distance=1.5cm and 2cm]

% Chapters
\node[chapter] (chap1) {第一章\\绪论};
\node[chapter, right=of chap1] (chap2) {第二章\\背景知识\\与相关技术};
\node[chapter, right=of chap2] (chap3) {第三章\\系统设计\\与实现};
\node[chapter, right=of chap3] (chap4) {第四章\\实验设计\\与结果分析};
\node[chapter, right=of chap4] (chap5) {第五章\\总结\\与展望};

% Sections for Chapter 1
\node[section, below=0.8cm of chap1, xshift=-1.2cm] (sec11) {研究背景\\与意义};
\node[section, below=0.6cm of sec11] (sec12) {国内外\\研究现状};
\node[section, below=0.6cm of sec12] (sec13) {研究内容\\与贡献};
\node[section, below=0.6cm of sec13] (sec14) {论文组织\\结构};

% Sections for Chapter 2
\node[section, below=0.8cm of chap2, xshift=-0.8cm] (sec21) {物联网设备\\硬件基础};
\node[section, below=0.6cm of sec21] (sec22) {物联网操作\\系统与驱动};
\node[section, below=0.6cm of sec22] (sec23) {大模型代码\\分析框架};
\node[section, below=0.6cm of sec23] (sec24) {静态分析\\技术};
\node[section, below=0.6cm of sec24] (sec25) {固件仿真与\\模糊测试};

% Sections for Chapter 3
\node[section, below=0.8cm of chap3, xshift=-0.8cm] (sec31) {系统整体\\架构设计};
\node[section, below=0.6cm of sec31] (sec32) {硬件依赖\\代码识别};
\node[section, below=0.6cm of sec32] (sec33) {驱动函数\\分析与替换};
\node[section, below=0.6cm of sec33] (sec34) {动态测试\\反馈闭环};

% Sections for Chapter 4
\node[section, below=0.8cm of chap4, xshift=-0.6cm] (sec41) {实验环境\\配置};
\node[section, below=0.6cm of sec41] (sec42) {实验设计\\方案};
\node[section, below=0.6cm of sec42] (sec43) {实验结果\\与分析};
\node[section, below=0.6cm of sec43] (sec44) {实验讨论\\与分析};

% Sections for Chapter 5
\node[section, below=0.8cm of chap5, xshift=-0.4cm] (sec51) {工作\\总结};
\node[section, below=0.6cm of sec51] (sec52) {不足与\\局限性};
\node[section, below=0.6cm of sec52] (sec53) {未来\\展望};

% Arrows connecting chapters
\draw[arrow] (chap1) -- (chap2);
\draw[arrow] (chap2) -- (chap3);
\draw[arrow] (chap3) -- (chap4);
\draw[arrow] (chap4) -- (chap5);

% Dashed lines from chapters to sections
\draw[dashed, gray, thick] (chap1) -- (sec11);
\draw[dashed, gray, thick] (chap1) -- (sec12);
\draw[dashed, gray, thick] (chap1) -- (sec13);
\draw[dashed, gray, thick] (chap1) -- (sec14);

\draw[dashed, gray, thick] (chap2) -- (sec21);
\draw[dashed, gray, thick] (chap2) -- (sec22);
\draw[dashed, gray, thick] (chap2) -- (sec23);
\draw[dashed, gray, thick] (chap2) -- (sec24);
\draw[dashed, gray, thick] (chap2) -- (sec25);

\draw[dashed, gray, thick] (chap3) -- (sec31);
\draw[dashed, gray, thick] (chap3) -- (sec32);
\draw[dashed, gray, thick] (chap3) -- (sec33);
\draw[dashed, gray, thick] (chap3) -- (sec34);

\draw[dashed, gray, thick] (chap4) -- (sec41);
\draw[dashed, gray, thick] (chap4) -- (sec42);
\draw[dashed, gray, thick] (chap4) -- (sec43);
\draw[dashed, gray, thick] (chap4) -- (sec44);

\draw[dashed, gray, thick] (chap5) -- (sec51);
\draw[dashed, gray, thick] (chap5) -- (sec52);
\draw[dashed, gray, thick] (chap5) -- (sec53);

% Title
\node[above=0.5cm of chap1, font=\large\bfseries, color=blue!80] (title) {论文组织结构};
\draw[dotted, blue!80, thick] (title) -- (chap1);

\end{tikzpicture}
\end{document}
